\begin{abstract}
传统辐射源定位通常先估计站间到达时间差，再求解源位置。当多个未知波形占用同一频带时，混合信号增加了时延提取和跨站关联的难度。直接位置估计（DPD）利用多站观测，在统一模型中联合估计源位置。然而，粗搜索中排名靠前的位置未必对应精修后的最佳多源解；逐源精修若固定信号系数，又可能因位置与波形估计的耦合而产生大量小步迭代。针对这两个求解问题，本文研究同步单天线站在已知自由空间相对幅度与基带时延模型下的未知宽带双源定位。首先以最小二乘消去未知波形，得到仅关于位置的优化目标；在此基础上，构造多候选联合精修交替投影（AP）与波形消元块NOMP。前者保留少量条件候选，分别联合调整两个源的位置，按最终拟合残差择优；后者在每次局部位置试探时重估信号系数，以消元后的曲率计算更新方向。条件投影和曲率推导给出两种策略的理论联系。在20个源站布局、六档功率差的2400条观测中，五候选联合AP将双源20\,m成功率由单候选的95\%提高至99\%，恢复96条观测且未损害原成功样本，平均求解时间为0.718\,s。相同停止精度下，单候选联合精修较交替精修减少约9.84\%的平均耗时。共同240条观测上的八类方法比较进一步展示了候选搜索与逐源拟合的可靠性--成本差异。
\end{abstract}
\noindent\textbf{关键词：}直接位置估计；强弱重叠源；交替投影；块NOMP；可分离非线性最小二乘
